%% ma: L11-B03-0011
%% lop: 11
%% chuong: 1
%% bai: 3
%% dang: 11.3.8
%% loai: DS
%% muc: [TH, NB, TH, TH]
%% dap_an: "SĐĐS"
%% nguon: "de-cuong-giua-ky-1-toan-11-nam-2026-2027-truong-thpt-hoang-van-thu-ha-noi-2-8.pdf (M08, Đề cương giữa kì 1 Toán 11 THPT Hoàng Văn Thụ, Hà Nội), Phần 2 Câu 2"
%% kien_thuc: "Bài 2 (công thức lượng giác), Bài 3 (hàm số lượng giác)"
%% trang_thai: cho-duyet
%% ly_do: "Dạng mới đề xuất 11.3.8"
%% ghi_chu: "Ý a: hiểu \"chu kì\" theo SGK là số dương nhỏ nhất T thỏa f(x + T) = f(x), nên chu kì là π và ý a sai."
\begin{ex}
Cho hàm số $y=f(x)=2\sin^2x-5$. Xét tính đúng sai của các khẳng định sau:
\choiceTF{Hàm số tuần hoàn với chu kì $2\pi$}{\True Hàm số là một hàm số chẵn}{\True Giá trị lớn nhất của hàm số đạt được khi $x=\dfrac{\pi}{2}+k\pi$, với $k\in\mathbb Z$}{Giá trị nhỏ nhất của hàm số bằng $-3$}
\loigiai{a) Sai: $f(x)=1-\cos2x-5=-4-\cos2x$ tuần hoàn với chu kì $\pi$.

b) Đúng: $f(-x)=2\sin^2(-x)-5=f(x)$ với mọi $x\in\mathbb R$.

c) Đúng: $f(x)\le2-5=-3$, dấu bằng khi $\sin^2x=1\Leftrightarrow x=\dfrac\pi2+k\pi$.

d) Sai: $f(x)\ge-5$, dấu bằng khi $\sin x=0$; giá trị nhỏ nhất là $-5$.}
\end{ex}
